% Source: https://artofproblemsolving.com/community/c3h2417951_inequality_of_the_triangle
% Stable identity: aops-topic-2417951-q01
Let the lengths be $a_1<a_2<\cdots<a_7$. Three lengths $x<y<z$ form a triangle exactly when $x+y>z$ (the other two triangle inequalities hold automatically, since $z$ is the largest).

Suppose that no three of the segments form a triangle. Then in particular, for $i=1,2,\dots,5$, the lengths $a_i<a_{i+1}<a_{i+2}$ do not, so
\[a_{i+2}\ge a_i+a_{i+1}.\]
Starting from $a_1\ge200$ and $a_2\ge200$, this gives
\[a_3\ge400,\quad a_4\ge600,\quad a_5\ge1000,\quad a_6\ge1600,\quad a_7\ge2600,\]
the lower bounds growing like the Fibonacci numbers $1,1,2,3,5,8,13$ times $200$:

\begin{asy}
usepackage("amsmath");
size(7cm);
defaultpen(fontsize(10pt));
real[] m={200,200,400,600,1000,1600,2600}; real s=0.002;
for(int i=0;i<7;++i){ real y=-0.45*i; fill(box((0,y-0.15),(s*m[i],y+0.15)),rgb(0.75,0.85,1)); draw(box((0,y-0.15),(s*m[i],y+0.15)));
 label("$a_"+string(i+1)+"\ge"+string(m[i])+"$",(0,y),W,fontsize(8pt)); }
draw((s*2009,0.4)--(s*2009,-3.05),red+dashed); label("$2009$",(s*2009,0.4),N,red);
\end{asy}

But $a_7\le2009$, a contradiction. Hence some three of the segments (in fact, three with consecutive lengths) form a triangle. $\blacksquare$
